% !PRINTED 383
\weberpage{383}
%<<CONT
\gesperrt{a) Every permutation $\pi$ of the roots of the double tangent equation which is applicable to all rational equations between these roots belongs to $P$.}

This follows thus: If $\pi$ is applicable to all rational equations between the roots, then the same holds of the powers of $\pi$. By 1., a syzygetic triple can never be carried by $\pi$ into an azygetic one, or conversely; for if $\xi_1\xi_2\xi_3$ is a syzygetic triple, $\pi$ is applicable to the equation (3); hence $\xi_1\xi_2\xi_3$ cannot pass into an azygetic triple. And the converse also is not possible, because otherwise a syzygetic triple would be carried into an azygetic one by $\pi^{-1}$. Hence any complete seven-system also passes, under $\pi$, again into such a system in some arrangement or other, and if one then applies $\pi$ to all equations of the form (5), there results a permutation of the $\xi_i, \xi_{ik}$ which belongs to $P$.

\gesperrt{b) Every rational equation between the roots of the double tangent equation admits all permutations of the group $P$.}

A rational relation between the roots has the form

\weq{6}{\Phi(\xi_1,\dots,\xi_7,\xi_{12},\dots,a\dots) = 0,}

in which $\Phi$ is a rational function, and $a\dots$ denote the ratios of the coefficients of $f$. Herein one can, by (5), express the $\xi_{12},\xi_{13},\dots$ rationally in terms of $\xi_1,\xi_2,\dots,\xi_7$, and by \S.\,100, II., since the curve $f$ is rationally determined by a complete seven-system of its double tangents, the $a$ can then be expressed rationally (with only numerical coefficients) in terms of the seven pairs of quantities (2)

\[
\xi_1, \Theta_1;\ \xi_2, \Theta_2;\ \dots;\ \xi_7, \Theta_7
\]

But since these can be given entirely arbitrarily, the equation (6) must pass into an identity under these substitutions. The equation (6) must therefore also remain correct if one replaces in it the $\xi_1,\xi_2,\dots,\xi_7$ by another seven-system, or also by the same one in another order, and at the same time carries out among the $\xi_{12},\dots$ the permutation prescribed by the formula (5), that is, if one performs among the roots of the double tangent equation a permutation of the group $P$.
